% !PRINTED 42
\weberpage{42}

%<<CONT
und wir zählen nun wieder ab, wie viel verschiedene Elemente
in dieser Form enthalten sind. Die beiden Werthe für $z$ müssen
dann übereinstimmen.

Zählen wir, wie oben in (7), die Anzahl der verschiedenen
in der Form
\weq{12}{B_1^{\beta_1 p_s m} B_2^{\beta_2 p_s m} \ldots B_{s-1}^{\beta_{s-1} p_s m}}
enthaltenen Elemente, so ergiebt sich auch hierfür nach der Annahme
 (6) der Werth $z$. Es kann daher in der Form
\weq{13}{B_s^{\beta_s p_s m} \ldots B_\mu^{\beta_\mu p_s m}}
nur das einzige Element $1$ enthalten sein, woraus zu schliessen
ist, dass $p_s$ durch $p_s'$ theilbar sein muss. Das ist aber mit (10)
nur unter der Voraussetzung vereinbar, dass
\weq{14}{p_s = p_s' \text{ ist.}}

Hiermit ist unser Theorem II. bewiesen.

Die in den Gradzahlen einer Basis von $S$ enthaltenen Primzahlpotenzen
 sind also von der besonderen Wahl der Basis ganz
unabhängig und wir nennen sie daher die \gesperrt{Invarianten der
Gruppe}. Das Product aller Invarianten ist gleich dem Grade $n$
der Gruppe.

Nach \S.\,9 können wir für $S$ eine Basis bestimmen, bei der
die Grade der Elemente lauter Primzahlpotenzen sind. Nach
dem Theorem II. sind diese Grade die Invarianten der Gruppe
und sind also durch die Gruppe vollständig bestimmt.

Dass die Invarianten auch die Natur der Gruppe vollständig
bestimmen, ergiebt sich aus dem Satze:

III. \gesperrt{Zwei Gruppen mit denselben Invarianten sind
isomorph, und isomorphe Gruppen haben dieselben
 Invarianten.}

Wenn nämlich zwei Gruppen $S$ und $S'$ der Anzahl und dem
Grade nach übereinstimmende Basiselemente haben, wenn etwa
\[ \Theta = A_1^{\alpha_1} A_2^{\alpha_2} \ldots A_\nu^{\alpha_\nu} \]
die Elemente von $S$ sind, und
\[ \Theta' = B_1^{\alpha_1} B_2^{\alpha_2} \ldots B_\nu^{\alpha_\nu} \]
die Elemente von $S'$, wo die $\alpha_1, \alpha_2 \ldots \alpha_\nu$ in beiden Fällen Restsysteme
 nach den Moduln $a_1, a_2 \ldots a_\nu$ durchlaufen, so brauchen
wir nur $\Theta$ und $\Theta'$ dann einander entsprechen zu lassen, wenn
%CONT>>
